\documentclass[a4paper,12pt]{article}

% ------------------------------------------------------------
% Кодировка, язык, математика
% ------------------------------------------------------------
\usepackage[T2A]{fontenc}
\usepackage[utf8]{inputenc}
\usepackage[russian]{babel}
\usepackage{amsmath,amssymb,amsthm,mathtools}
\usepackage{helvet}
\renewcommand{\familydefault}{\sfdefault}
\usepackage{icomma}

% ------------------------------------------------------------
% Вёрстка и типографика
% ------------------------------------------------------------
\usepackage[
  left=18mm,
  right=18mm,
  top=18mm,
  bottom=20mm
]{geometry}
\usepackage{microtype}
\usepackage{parskip}
\usepackage{setspace}
\onehalfspacing
\usepackage{enumitem}
\usepackage{tabularx,booktabs,longtable,array}
\usepackage{xcolor}
\usepackage{fancyhdr}
\usepackage{hyperref}
\usepackage[most]{tcolorbox}

% ------------------------------------------------------------
% Графика
% ------------------------------------------------------------
\usepackage{tikz}
\usetikzlibrary{calc}
\usetikzlibrary{arrows.meta,positioning,shapes.geometric}
\usepackage{tkz-euclide}
\usepackage{pgfplots}
\pgfplotsset{compat=1.18}

% ------------------------------------------------------------
% Цветовая палитра
% ------------------------------------------------------------
\definecolor{accent}{HTML}{008080}
\definecolor{darkaccent}{HTML}{004D4D}
\definecolor{textgray}{HTML}{333333}
\definecolor{softaccent}{HTML}{EAF4F4}
\definecolor{softgray}{HTML}{F4F5F5}
\definecolor{warning}{HTML}{8C5A3C}

\color{textgray}

% ------------------------------------------------------------
% Служебные команды
% ------------------------------------------------------------
\newcommand{\topic}{Геометрический смысл производной: графики функции и производной}

\newcommand{\tg}{\mathop{\mathrm{tg}}\nolimits}
\newcommand{\ctg}{\mathop{\mathrm{ctg}}\nolimits}
\newcommand{\arctg}{\mathop{\mathrm{arctg}}\nolimits}

\hypersetup{
  colorlinks=true,
  linkcolor=darkaccent,
  urlcolor=darkaccent,
  pdftitle={Чек-лист: геометрический смысл производной},
  pdfauthor={Лёвин Артём Александрович}
}

\pagestyle{fancy}
\fancyhf{}
\fancyhead[L]{\small\color{darkaccent}\topic}
\fancyhead[R]{\small\color{textgray}\student}
\fancyfoot[C]{\small\thepage}
\renewcommand{\headrulewidth}{0.3pt}
\renewcommand{\headrule}{\hbox to\headwidth{\color{accent}\leaders\hrule height \headrulewidth\hfill}}

% ------------------------------------------------------------
% Блоки
% ------------------------------------------------------------
\newtcolorbox{keybox}[1][]{
  enhanced,
  breakable,
  colback=softaccent,
  colframe=accent,
  boxrule=0.6pt,
  arc=1.5mm,
  left=3mm,
  right=3mm,
  top=2.5mm,
  bottom=2.5mm,
  title=#1,
  fonttitle=\bfseries,
  coltitle=darkaccent
}

\newtcolorbox{noticebox}[1][]{
  enhanced,
  breakable,
  colback=softgray,
  colframe=black!25,
  boxrule=0.5pt,
  arc=1.5mm,
  left=3mm,
  right=3mm,
  top=2.5mm,
  bottom=2.5mm,
  title=#1,
  fonttitle=\bfseries,
  coltitle=textgray
}

\newtcolorbox{errorbox}[1][]{
  enhanced,
  breakable,
  colback=white,
  colframe=warning,
  boxrule=0.6pt,
  arc=1.5mm,
  left=3mm,
  right=3mm,
  top=2.5mm,
  bottom=2.5mm,
  title=#1,
  fonttitle=\bfseries,
  coltitle=warning
}

% ------------------------------------------------------------
% TikZ-стили для блок-схемы
% ------------------------------------------------------------
\tikzset{
  flowstart/.style={
    ellipse,
    draw=accent,
    very thick,
    align=center,
    minimum width=42mm,
    minimum height=11mm,
    fill=softaccent,
    font=\small
  },
  flowaction/.style={
    rounded rectangle,
    draw=accent,
    thick,
    align=center,
    text width=72mm,
    minimum height=12mm,
    fill=white,
    font=\small
  },
  flowdecision/.style={
    diamond,
    aspect=2.2,
    draw=darkaccent,
    thick,
    align=center,
    text width=42mm,
    inner sep=1.5mm,
    fill=softgray,
    font=\small
  },
  flowresult/.style={
    rounded rectangle,
    draw=darkaccent,
    very thick,
    align=center,
    text width=70mm,
    minimum height=12mm,
    fill=softaccent,
    font=\small
  },
  flowarrow/.style={
    -{Stealth[length=3mm,width=2mm]},
    thick,
    draw=darkaccent
  }
}

\begin{document}

% ============================================================
% ТИТУЛЬНЫЙ ЛИСТ
% ============================================================
\begin{titlepage}
\thispagestyle{empty}

\begin{center}

\vspace*{20mm}

{\Large\bfseries\color{darkaccent} Чек-лист}

\vspace{8mm}

{\huge\bfseries
Геометрический смысл производной}

\vspace{5mm}

{\Large
Графики функции и производной}

\vspace{20mm}

{\large
ЕГЭ по профильной математике}

\vspace{13mm}

{\large
\textbf{Ученица:} Николь Сарикисьянц}

\vspace{6mm}

{\large
\textbf{Преподаватель:} Лёвин Артём Александрович}

\vspace{6mm}

{\large
\today}

\vfill

{\normalsize
\teacher\\[2mm]
эксклюзивно для \student}

\vspace{18mm}

\begin{tikzpicture}[remember picture,overlay]
  \draw[
    accent,
    line width=1.1pt,
    opacity=0.65
  ]
  ($(current page.south west)+(20mm,17mm)$)
  .. controls
  ($(current page.south west)+(65mm,28mm)$)
  and
  ($(current page.south east)+(-70mm,6mm)$)
  ..
  ($(current page.south east)+(-20mm,18mm)$);
\end{tikzpicture}

\end{center}
\end{titlepage}

% ============================================================
% ОСНОВНОЙ ЧЕК-ЛИСТ
% ============================================================
\section*{1. Главная идея занятия}

Производная связывает \textbf{поведение функции} с \textbf{наклоном её графика}.
В задачах ЕГЭ прежде всего необходимо определить, чей именно график изображён:
\[
y=f(x)
\qquad\text{или}\qquad
y=f'(x).
\]

От этого зависит весь дальнейший способ решения.

\begin{keybox}[Базовая связь]
Для дифференцируемой функции:
\[
f'(x)>0
\quad\Longrightarrow\quad
f(x)\text{ возрастает},
\]
\[
f'(x)<0
\quad\Longrightarrow\quad
f(x)\text{ убывает}.
\]

Если
\[
f'(x)=0,
\]
то касательная к графику функции в этой точке горизонтальна.
\end{keybox}

\section*{2. Первый вопрос: чей график дан?}

\begin{center}
\renewcommand{\arraystretch}{1.45}
\begin{tabularx}{\linewidth}{>{\bfseries}p{0.27\linewidth} X X}
\toprule
Что изображено & Что можно считывать непосредственно & Что приходится восстанавливать \\
\midrule
$y=f(x)$
&
значения функции, возрастание и убывание, экстремумы, горизонтальные касательные
&
знак и отдельные значения производной
\\
$y=f'(x)$
&
знак производной, нули производной
&
возрастание и убывание функции, характер некоторых экстремумов
\\
\bottomrule
\end{tabularx}
\end{center}

\begin{noticebox}[Контрольный вопрос]
Перед любыми вычислениями сформулируйте словами:

\begin{center}
\textbf{«Я смотрю на график функции или на график производной?»}
\end{center}
\end{noticebox}

% ============================================================
\section*{3. Если дан график функции \(y=f(x)\)}

\subsection*{3.1. Как определить знак производной}

Смотрите на направление движения графика функции слева направо.

\[
\boxed{\text{график идёт вверх}}
\quad\Longrightarrow\quad
f'(x)>0,
\]

\[
\boxed{\text{график идёт вниз}}
\quad\Longrightarrow\quad
f'(x)<0.
\]

Следовательно, если требуется найти целые значения \(x\), при которых
\[
f'(x)<0,
\]
необходимо выделить участки \textbf{убывания функции} и выбрать подходящие
целые абсциссы.

\begin{keybox}[Важно]
Все значения считываются \textbf{по оси \(x\)}.

Ордината точки графика в задачах вида
\[
f'(x)>0,\qquad f'(x)<0
\]
обычно не является искомым числом.
\end{keybox}

\subsection*{3.2. Когда \(f'(x)=0\)}

Условие
\[
f'(x_0)=0
\]
означает, что в точке с абсциссой \(x_0\) касательная горизонтальна.

Это возможно:

\begin{itemize}[leftmargin=7mm]
  \item в точке локального максимума;
  \item в точке локального минимума;
  \item в стационарной точке, которая экстремумом не является.
\end{itemize}

\begin{errorbox}[Типичная ошибка]
Нельзя считать, что
\[
f'(x_0)=0
\]
автоматически означает наличие экстремума.

Например, у функции
\[
f(x)=x^3
\]
имеем
\[
f'(0)=0,
\]
хотя в точке \(x=0\) функция продолжает возрастать и экстремума нет.
\end{errorbox}

\subsection*{3.3. Экстремум функции}

Если функция дифференцируема и в точке внутреннего экстремума существует производная, то
\[
f'(x_0)=0.
\]

Для подтверждения экстремума особенно полезно следить за изменением знака производной:

\[
+\ \longrightarrow\ -
\quad\Rightarrow\quad
\text{локальный максимум},
\]

\[
-\ \longrightarrow\ +
\quad\Rightarrow\quad
\text{локальный минимум}.
\]

Если знак производной при прохождении через ноль не меняется, экстремума может не быть.

% ============================================================
\section*{4. Если дан график производной \(y=f'(x)\)}

Это ключевой тип задач занятия.

Поскольку
\[
y=f'(x),
\]
знак ординаты точки графика одновременно является знаком производной.

\subsection*{4.1. График выше оси \(Ox\)}

Если график производной расположен выше оси \(Ox\), то
\[
f'(x)>0.
\]

Следовательно,
\[
f(x)\text{ возрастает}.
\]

\subsection*{4.2. График ниже оси \(Ox\)}

Если график производной расположен ниже оси \(Ox\), то
\[
f'(x)<0.
\]

Следовательно,
\[
f(x)\text{ убывает}.
\]

\subsection*{4.3. Пересечение с осью \(Ox\)}

В точках пересечения графика производной с осью абсцисс
\[
y=0.
\]

Так как
\[
y=f'(x),
\]
получаем
\[
f'(x)=0.
\]

\begin{keybox}[Сверхкороткая памятка]
Для графика \(y=f'(x)\):
\[
\text{выше }Ox \Rightarrow f\uparrow,
\]
\[
\text{ниже }Ox \Rightarrow f\downarrow,
\]
\[
\text{на }Ox \Rightarrow f'=0.
\]
\end{keybox}

% ============================================================
\section*{5. Максимум и минимум по графику производной}

Если производная меняет знак

\[
f'(x):
\quad
+\ \longrightarrow\ -,
\]
то функция сначала возрастает, затем убывает.

Следовательно, в соответствующей точке функция имеет локальный максимум.

Если
\[
f'(x):
\quad
-\ \longrightarrow\ +,
\]
то функция сначала убывает, затем возрастает.

Следовательно, функция имеет локальный минимум.

\begin{center}
\renewcommand{\arraystretch}{1.5}
\begin{tabularx}{0.92\linewidth}{>{\centering\arraybackslash}p{0.27\linewidth}
>{\centering\arraybackslash}p{0.27\linewidth}
>{\centering\arraybackslash}X}
\toprule
Знак \(f'\) слева & Знак \(f'\) справа & Что происходит с \(f\) \\
\midrule
\( + \) & \( - \) & максимум \\
\( - \) & \( + \) & минимум \\
\( + \) & \( + \) & экстремум не следует \\
\( - \) & \( - \) & экстремум не следует \\
\bottomrule
\end{tabularx}
\end{center}

% ============================================================
\section*{6. Наибольшее и наименьшее значение на отрезке}

В задачах может быть дан график производной и указан отрезок
\[
[a;b].
\]

Здесь требуется анализировать только поведение функции внутри указанного промежутка.

\subsection*{6.1. Производная положительна на всём отрезке}

Если
\[
f'(x)>0
\qquad
\text{для всех }x\in(a;b),
\]
то функция возрастает на \([a;b]\).

Тогда
\[
\min_{[a;b]}f(x)=f(a),
\qquad
\max_{[a;b]}f(x)=f(b).
\]

Следовательно, если спрашивают значение \(x\), при котором достигается минимум,
ответ:
\[
x=a.
\]

Если спрашивают значение \(x\), при котором достигается максимум,
ответ:
\[
x=b.
\]

\subsection*{6.2. Производная отрицательна на всём отрезке}

Если
\[
f'(x)<0
\qquad
\text{для всех }x\in(a;b),
\]
то функция убывает.

Поэтому
\[
\max_{[a;b]}f(x)=f(a),
\qquad
\min_{[a;b]}f(x)=f(b).
\]

\begin{noticebox}[Очень важная мысль]
Для определения точки максимума или минимума на промежутке часто
\textbf{не требуется знать точные значения \(f(x)\)}.

Достаточно установить, возрастает функция или убывает.
\end{noticebox}

\subsection*{6.3. Почему эскиз можно рисовать произвольно по вертикали}

По графику производной обычно нельзя восстановить абсолютное положение графика функции по оси \(Oy\) без дополнительной информации.

Например, если известно лишь
\[
f'(x)>0,
\]
можно утверждать, что функция возрастает, однако её конкретные значения могут быть различными.

Поэтому схематичные возрастающие линии
\[
f_1(x),\quad f_2(x),\quad f_3(x)
\]
могут располагаться на разной высоте, сохраняя один и тот же характер монотонности.

% ============================================================
\section*{7. Касательная, параллельная заданной прямой}

Геометрический смысл производной:
\[
f'(x_0)=k,
\]
где \(k\) --- угловой коэффициент касательной к графику функции в точке с абсциссой \(x_0\).

Если касательная параллельна прямой
\[
y=kx+b,
\]
то её угловой коэффициент также равен \(k\), поэтому
\[
f'(x_0)=k.
\]

\begin{keybox}[Частный случай]
Если касательная параллельна горизонтальной прямой
\[
y=c,
\]
то её угловой коэффициент равен нулю.

Следовательно,
\[
f'(x_0)=0.
\]

Например, для прямой
\[
y=-3
\]
получаем условие
\[
f'(x_0)=0.
\]
\end{keybox}

\begin{errorbox}[Типичная ошибка]
Правило
\[
\text{«касательная параллельна прямой»}
\Rightarrow f'(x)=0
\]
верно только для \textbf{горизонтальной} прямой вида
\[
y=c.
\]

Для прямой
\[
y=5x-2
\]
необходимо искать точки, где
\[
f'(x)=5.
\]
\end{errorbox}

% ============================================================
\section*{8. Формулировки ЕГЭ: что именно записывать в ответ}

После математического анализа обязательно перечитайте вопрос.

Могут попросить:

\begin{itemize}[leftmargin=7mm]
  \item количество подходящих точек;
  \item конкретное значение \(x\);
  \item сумму целых значений \(x\);
  \item длину промежутка;
  \item длину наибольшего промежутка;
  \item точку максимума;
  \item точку минимума;
  \item значение \(x\), при котором функция принимает наибольшее значение;
  \item значение \(x\), при котором функция принимает наименьшее значение.
\end{itemize}

\begin{keybox}[Финальная проверка]
Даже правильно найденные промежутки ещё не гарантируют правильный ответ.

Последний шаг решения:
\[
\boxed{\text{выполнить именно то действие, которое требует условие}}
\]
--- посчитать, сложить, измерить длину или выбрать нужную абсциссу.
\end{keybox}

% ============================================================
\section*{9. Работа с целыми значениями \(x\)}

Если требуется найти целые значения \(x\), принадлежащие нужным промежуткам, смотрите только на ось абсцисс.

Например, если условие выполняется при
\[
-5{,}7<x<-2{,}4,
\]
то целые значения:
\[
-5,\,-4,\,-3.
\]

Их количество:
\[
3.
\]

Их сумма:
\[
-5-4-3=-12.
\]

\begin{errorbox}[Граница промежутка]
Если требуется строгое неравенство
\[
f'(x)<0,
\]
точка, где
\[
f'(x)=0,
\]
не включается.

При условии
\[
f'(x)\le 0
\]
такая точка, напротив, может входить в ответ.
\end{errorbox}

% ============================================================
\section*{10. Работа с заданным промежутком}

Даже если график изображён на большом диапазоне, анализировать необходимо именно тот промежуток, который указан в задаче.

Если требуется исследовать
\[
x\in[-4;1],
\]
участки графика левее \(-4\) и правее \(1\) в решение не входят.

\begin{noticebox}[Проверка перед ответом]
Отметьте границы нужного промежутка непосредственно на рисунке.

Это снижает риск включить лишнюю точку или пропустить нужную.
\end{noticebox}

% ============================================================
\section*{11. Универсальные четыре шага}

Для большинства задач этого типа удобно использовать один шаблон.

\begin{enumerate}[leftmargin=8mm,label=\textbf{\arabic*.}]
  \item Определите, чей график дан:
  \[
  y=f(x)\quad\text{или}\quad y=f'(x).
  \]

  \item Переведите условие задачи на язык производной:
  \[
  f'(x)>0,\quad
  f'(x)<0,\quad
  f'(x)=0
  \]
  либо установите требуемое поведение функции.

  \item Найдите нужные точки или промежутки на графике.

  \item Перечитайте вопрос и сформируйте ответ в требуемом виде:
  количество, сумма, длина, абсцисса максимума или минимума и т.\,д.
\end{enumerate}


% ============================================================
\section*{12. Мини-чек-лист перед сдачей задачи}

Перед записью ответа проверьте:

\begin{itemize}[leftmargin=7mm,label=\(\square\)]
  \item Я определила, изображён график \(f\) или \(f'\).
  \item Я выделила именно тот промежуток, который указан в условии.
  \item Я правильно связала знак \(f'\) с возрастанием и убыванием \(f\).
  \item При условии \(f'(x)=0\) я ищу нули производной, а не автоматически все экстремумы.
  \item Если требуется экстремум, я проверила изменение знака производной.
  \item Если дана касательная, я нашла её угловой коэффициент.
  \item Целые значения я считываю по оси \(x\).
  \item Точки с \(f'(x)=0\) не включаю в условие \(f'(x)<0\) или \(f'(x)>0\).
  \item Я перечитала формулировку: требуется количество, сумма, длина или конкретное значение \(x\).
  \item Ответ записан в требуемом формате.
\end{itemize}

% ============================================================
\section*{13. Типичные ошибки}

\begin{enumerate}[leftmargin=8mm]
  \item \textbf{Перепутать графики \(f\) и \(f'\).}

  Если изображён \(f'\), участок ниже оси \(Ox\) означает убывание
  исходной функции.

  \item \textbf{Считать ординаты вместо абсцисс.}

  В большинстве задач данного типа требуется значение \(x\).

  \item \textbf{Включить нуль производной в строгое неравенство.}

  При
  \[
  f'(x)<0
  \]
  точки с
  \[
  f'(x)=0
  \]
  исключаются.

  \item \textbf{Считать каждый нуль \(f'\) экстремумом.}

  Экстремум требует соответствующего изменения поведения функции.

  \item \textbf{Забыть о границах заданного промежутка.}

  В задаче анализируется только заявленная область.

  \item \textbf{Найти точки правильно, но неверно оформить ответ.}

  Например, условие требует сумму целых значений, а записано их количество.

  \item \textbf{Принять любую параллельную прямую за горизонтальную.}

  Для
  \[
  y=kx+b
  \]
  требуется
  \[
  f'(x)=k.
  \]
\end{enumerate}

% ============================================================
\section*{14. Сводная таблица}

\begin{center}
\renewcommand{\arraystretch}{1.45}
\begin{tabularx}{\linewidth}{
>{\raggedright\arraybackslash}p{0.27\linewidth}
>{\centering\arraybackslash}p{0.22\linewidth}
X}
\toprule
Наблюдение & Запись & Вывод \\
\midrule

Функция возрастает
&
\(f'(x)>0\)
&
производная положительна
\\

Функция убывает
&
\(f'(x)<0\)
&
производная отрицательна
\\

Горизонтальная касательная
&
\(f'(x)=0\)
&
стационарная точка
\\

\(f'\): \(+\to-\)
&
\(f'(x_0)=0\)
&
локальный максимум \(f\)
\\

\(f'\): \(-\to+\)
&
\(f'(x_0)=0\)
&
локальный минимум \(f\)
\\

График \(f'\) выше \(Ox\)
&
\(f'(x)>0\)
&
\(f\) возрастает
\\

График \(f'\) ниже \(Ox\)
&
\(f'(x)<0\)
&
\(f\) убывает
\\

График \(f'\) пересекает \(Ox\)
&
\(f'(x)=0\)
&
кандидат на стационарную точку
\\

Касательная параллельна \(y=kx+b\)
&
\(f'(x)=k\)
&
одинаковый угловой коэффициент
\\

\bottomrule
\end{tabularx}
\end{center}

\end{document}